\section{Estratégia de operação do manipulador robótico}
\label{sec:estrategia_operacao_MR}

A operação do \gls{mr} é organizada de forma coerente com as fases de aproximação relativa em \gls{lvlh} e com os critérios de atitude da espaçonave.
Em termos operacionais, o manipulador permanece inicialmente travado (modo não-operacional), de modo que o conjunto base--\gls{mr} possa ser tratado como um corpo rígido durante a transferência orbital, o \textit{drift} em órbita de fase e a maior parte da aproximação curta.
Nessa fase, somente os atuadores da base são utilizados para regular posição e atitude em relação ao alvo, enquanto o \gls{mr} é mantido em configuração de espera.

Quando o perseguidor inicia a aproximação final, o \gls{mr} passa ao modo operacional.
A partir desse ponto, o manipulador é comandado a realizar o alinhamento fino do efetuador com o eixo da porta de acoplamento.
Em todas essas etapas, o movimento das juntas é planejado para respeitar limites de curso, velocidade e torque, de modo a evitar colisões com a base ou com o alvo.

A estratégia de operação leva em conta que qualquer movimentação do \gls{mr} gera reações na base, que são tratadas na dinâmica acoplada e devem ser compensadas pelo controle de atitude.
Assim, o planejamento das trajetórias articulares privilegia movimentos com variações moderadas de velocidade e aceleração.